% Generated draft: requires mathematical and source review.
\documentclass[11pt]{article}
\usepackage[margin=25mm]{geometry}
\usepackage{fontspec}
\setmainfont{Latin Modern Roman}
\usepackage{amsmath,amssymb,mathtools,unicode-math}
\setmathfont{Latin Modern Math}
\usepackage{enumitem,needspace,xurl,hyperref}
\setlength{\parindent}{0pt}
\setlength{\parskip}{5pt}
\setlength{\emergencystretch}{3em}
\newcommand{\N}{\mathbb{N}}
\newcommand{\Z}{\mathbb{Z}}
\newcommand{\Q}{\mathbb{Q}}
\newcommand{\R}{\mathbb{R}}
\newcommand{\C}{\mathbb{C}}
\newcommand{\F}{\mathbb{F}}
\newcommand{\A}{\mathbb{A}}
\newcommand{\PP}{\mathbb{P}}
\newcommand{\E}{\mathbb{E}}
\newcommand{\eps}{\varepsilon}
\newcommand{\dd}{\,\mathrm{d}}
\begin{document}
\section*{UnsolvedMath ID: 800781}
\textbf{Source problem code:} catalogue entry
\subsection{Definitions and mathematical statement}
Fix a positive integer \(n\) and use local coordinates centered at \(O\in\C^n\). A holomorphic local dynamical system at \(O\) is a germ \(f:(\C^n,O)\to(\C^n,O)\) of a holomorphic map fixing \(O\). Equivalently, it is an element of
\[
\mathcal H_n=\C_0\{z_1,\ldots,z_n\}^{\,n},
\]
the space of \(n\)-tuples of convergent power series with zero constant terms. No invertibility assumption on \(f\) is imposed.

For a chosen regularity \(r\), two such germs \(f_1,f_2\) are locally conjugate with regularity \(r\) if there is an invertible germ \(h:(\C^n,O)\to(\C^n,O)\), fixing \(O\), such that
\[
h\circ f_1=f_2\circ h.
\]
For holomorphic conjugacy, \(h\) is biholomorphic; for smooth conjugacy, \(h\) and \(h^{-1}\) are \(C^\infty\) as maps of the underlying real spaces; for \(C^k\) conjugacy, both are \(C^k\); and for topological conjugacy, \(h\) is a homeomorphism. The source does not specify the range of \(k\); the usual finite differentiability interpretation takes a specified positive integer \(k\).

\textbf{Explicit invariant problem in the supplied excerpt.} For each chosen regularity among holomorphic, smooth, \(C^k\), and topological, find a way to assign to every \(f\in\mathcal H_n\) a collection of objects \(I_r(f)\), called invariants, such that, for all \(f_1,f_2\in\mathcal H_n\),
\[
f_1\text{ and }f_2\text{ are locally conjugate with regularity }r
\quad\Longrightarrow\quad
I_r(f_1)=I_r(f_2).
\]
The source allows a possibly small collection of possibly computable objects, without imposing a size bound or a computability requirement. Such invariants are called \emph{complete} when, for all \(f_1,f_2\in\mathcal H_n\),
\[
f_1\text{ and }f_2\text{ are locally conjugate with regularity }r
\quad\Longleftrightarrow\quad
I_r(f_1)=I_r(f_2).
\]
The excerpt defines completeness but does not explicitly require that the requested invariants be complete.

\textbf{Supplied formal-conjugacy background.} Let
\[
\mathcal F_n=\C_0[[z_1,\ldots,z_n]]^{\,n}
\]
be the space of \(n\)-tuples of formal power series with zero constant terms. This space is closed under composition. An element \(\Phi\in\mathcal F_n\) has a compositional inverse in \(\mathcal F_n\) exactly when its linear part is an invertible linear map of \(\C^n\). Two elements \(f_1,f_2\in\mathcal F_n\) are \emph{formally conjugate} if there is such an invertible \(\Phi\) satisfying
\[
f_1=\Phi^{-1}\circ f_2\circ\Phi
\]
as an identity of formal power series. For holomorphic germs, formal conjugacy refers to their Taylor series. The supplied background states that holomorphic local conjugacy implies both formal and topological local conjugacy. It also reports examples of topologically conjugate holomorphic germs that are neither formally nor holomorphically conjugate, and examples of formally conjugate holomorphic germs that are neither holomorphically nor topologically conjugate.

\textbf{Unresolved scope qualification.} The record identifies the source item as Problem Q4, but the excerpt ends by introducing a final general question after the formal-conjugacy discussion, without providing that question. Thus the invariant problem above is the explicit task present in the supplied text; its identification with the complete intended Q4 formulation cannot be established from this record.
\subsection{Short English statement}
For holomorphic map germs fixing a point of \(\C^n\), seek invariants preserved by local conjugacy in each specified regularity category. Complete invariants would characterize conjugacy in both directions. The supplied excerpt also defines formal conjugacy, but omits the final question it introduces, leaving the precise intended scope of Problem Q4 unresolved.
\subsection{Sources}
\begin{itemize}
\item Local UnsolvedMath registry, supplied complete record with permanent ID \(800781\), exact source problem code \texttt{catalogue entry}, titled \emph{Local Discrete Holomorphic Dynamics --- Open Problem Q4}. The registry attributes the problem to Abate, with proposed year 2011, and labels it \emph{open}; this is supplied metadata, not an independently verified status assessment.
\item Abate, \emph{Open Problems in Local Discrete Holomorphic Dynamics} (2011), source item Problem Q4, \url{https://arxiv.org/abs/1107.4876}. The record specifies extraction by \texttt{arxiv\_latex} and an access date of 2026-10-16. This access date and attribution are supplied provenance and have not been independently verified here. The supplied excerpt ends before the final question introduced by its concluding sentence.
\end{itemize}
\end{document}
